% Restored from ca5273cd335633e0aea2e4930c897d732b8260de; current validation pending.
% Historical acceptance badges removed until the full dependency closure passes.
\subsection{The original-spin quantum-double Hamiltonian}
\label{sec:peps_qd_original_spins}

Let $G$ be a finite group and let $w,h\ge3$.
On the original lattice
$\Lambda=(\mathbb Z/2w\mathbb Z)\times(\mathbb Z/2h\mathbb Z)$,
place one copy of $\mathbb C[G]$ at every vertex.
Use the four-site tiling $t(v,i)=2v+b_i$ of
Definition~\ref{def:peps_tc_original_faces}, with
$b_0=(1,1)$, $b_1=(1,0)$, $b_2=(0,0)$, and $b_3=(0,1)$.
The tile interiors are the clockwise A faces in the arrow convention of
\cite[Section~7.2]{Schuch2010PEPS}; the intervening A faces have the
opposite geometric orientation. The face bijection in
Theorem~\ref{thm:peps_tc_face_geometry} is independent of $G$.
The period restriction is recorded in
\cite{gap:rmp_peps_examples_small_torus}.

\begin{lemma}[Products of commuting orthogonal projections]
    \label{lem:peps_qd_commuting_projection_product}
    \lean{Matrix.isStarProjection_list_prod,Matrix.list_prod_mulVec_eq_self_iff}
    Let $(P_j)_{j\in J}$ be pairwise commuting orthogonal projections on a
    finite-dimensional complex vector space. For a finite list
    $j_1,\ldots,j_n$, the product $Q=P_{j_1}\cdots P_{j_n}$ is an
    orthogonal projection, and
    \[
        Q\psi=\psi\quad\Longleftrightarrow\quad
        P_{j_k}\psi=\psi\text{ for every }1\le k\le n.
    \]
    The list may repeat indices, and its empty product is the identity.
\end{lemma}
\begin{proof}
    Commutation and idempotence give $P_{j_k}Q=Q$ for every listed factor.
    Applying this equality to a vector fixed by $Q$ proves one implication;
    successive application of the factors proves the converse.
    The product is self-adjoint and idempotent by the same commutation identities.
\end{proof}

\begin{definition}[Ordered products and original spin actions]
    \label{def:peps_qd_original_actions}
    \lean{TNLean.PEPS.quantumDoubleOriginalAHolonomy,
        TNLean.PEPS.quantumDoubleOriginalBOutgoing,
        TNLean.PEPS.quantumDoubleOriginalBIncoming,
        TNLean.PEPS.quantumDoubleOriginalBAction,
        TNLean.PEPS.quantumDoubleOriginalBPermutation}

    Write $\sigma_{v,i}=\sigma(t(v,i))$ for $\sigma\in G^\Lambda$,
    and let $\mathrm e,\mathrm n$ denote the right and upper coarse steps.
    The two A products are
    \begin{align}
        a_{\mathrm{in},v}(\sigma)
          &=\sigma_{v,0}\sigma_{v,1}\sigma_{v,2}\sigma_{v,3},\\
        a_{\mathrm{hole},v}(\sigma)
          &=\sigma_{v+\mathrm n,1}\sigma_{v,0}
            \sigma_{v+\mathrm e,3}\sigma_{v+\mathrm e+\mathrm n,2}.
        \label{eq:peps_qd_original_a}
    \end{align}
    For a B face $e$, denote its outgoing and incoming spin sets by
    $O_e$ and $I_e$. Across a right boundary and an upper boundary,
    respectively, these are
    \begin{align}
        O_{\mathrm e,v}&=\{t(v,1),t(v+\mathrm e,3)\},&
        I_{\mathrm e,v}&=\{t(v,0),t(v+\mathrm e,2)\},\\
        O_{\mathrm n,v}&=\{t(v,0),t(v+\mathrm n,2)\},&
        I_{\mathrm n,v}&=\{t(v,3),t(v+\mathrm n,1)\}.
        \label{eq:peps_qd_original_incidence}
    \end{align}
    For $u\in G$, the original four-spin operation is
    \begin{align}
        (V_e(u)\sigma)_p=
        \begin{cases}
            u\sigma_p,&p\in O_e,\\
            \sigma_pu^{-1},&p\in I_e,\\
            \sigma_p,&p\notin O_e\cup I_e.
        \end{cases}
        \label{eq:peps_qd_original_action}
    \end{align}
    Thus an outgoing arrow receives $L_u$ and an incoming arrow receives
    $R_u$, as in \cite[Section~7.2]{Schuch2010PEPS}.
\end{definition}

\begin{theorem}[Four-spin support and the group action]
    \label{thm:peps_qd_original_action_support}
    \lean{TNLean.PEPS.quantumDoubleOriginalBOutgoing_right,
        TNLean.PEPS.quantumDoubleOriginalBOutgoing_up,
        TNLean.PEPS.quantumDoubleOriginalBIncoming_right,
        TNLean.PEPS.quantumDoubleOriginalBIncoming_up,
        TNLean.PEPS.quantumDoubleOriginalB_incidence,
        TNLean.PEPS.quantumDoubleOriginalBAction_spectator,
        TNLean.PEPS.quantumDoubleOriginalBAction_one,
        TNLean.PEPS.quantumDoubleOriginalBAction_mul}

    \uses{def:peps_qd_original_actions,thm:peps_tc_face_geometry}
    The sets $O_e$ and $I_e$ are disjoint, and their union is exactly
    the four-site B face. Every other spin is fixed.
    Moreover, $V_e(1)=I$ and $V_e(uv)=V_e(u)V_e(v)$ on the whole
    original configuration space, including across either periodic seam.
\end{theorem}
\begin{proof}
    The four corner labels in~(\ref{eq:peps_qd_original_incidence}) are
    distinct under the bijective tiling. On outgoing spins, composition
    gives $u(vg)=(uv)g$; on incoming spins it gives
    $(gv^{-1})u^{-1}=g(uv)^{-1}$. Other spins remain fixed.
\end{proof}

\begin{definition}[Literal original A and B penalties]
    \label{def:peps_qd_original_hamiltonian}
    \lean{TNLean.PEPS.quantumDoubleOriginalATerm,
        TNLean.PEPS.quantumDoubleOriginalBMatrix,
        TNLean.PEPS.quantumDoubleOriginalBAverage,
        TNLean.PEPS.quantumDoubleOriginalBTerm,
        TNLean.PEPS.quantumDoubleOriginalTerm,
        TNLean.PEPS.quantumDoubleOriginalHamiltonian}

    \uses{def:peps_qd_original_actions}
    Let $M_e(u)\ket\sigma=\ket{V_e(u)\sigma}$. Define
    \begin{align}
        h_A^j&=I-\sum_{\sigma:a_j(\sigma)=1}\ket\sigma\bra\sigma,\\
        P_B^e&=\frac1{|G|}\sum_{u\in G}M_e(u),& h_B^e&=I-P_B^e,\\
        H&=\sum_j h_A^j+\sum_e h_B^e.
        \label{eq:peps_qd_original_hamiltonian}
    \end{align}
    Every spectator spin carries the identity. The factor $|G|^{-1}$
    corrects the unnormalized B sum printed in
    \cite[Section~7.2]{Schuch2010PEPS}; see
    \cite{gap:scp10_quantum_double_local_hamiltonian}.
\end{definition}

\begin{theorem}[Coefficientwise physical regrouping]
    \label{thm:peps_qd_original_blocking}
    \lean{TNLean.PEPS.quantumDoubleOriginalAHolonomy_block,
        TNLean.PEPS.quantumDoubleOriginalBAction_block,
        TNLean.PEPS.quantumDoubleOriginalBMatrix_apply,
        TNLean.PEPS.quantumDoubleUnblockOperator_apply,
        TNLean.PEPS.quantumDoubleOriginalATerm_eq_unblocked,
        TNLean.PEPS.quantumDoubleOriginalBMatrix_eq_unblocked,
        TNLean.PEPS.quantumDoubleOriginalBAverage_eq_unblocked,
        TNLean.PEPS.quantumDoubleOriginalBTerm_eq_unblocked,
        TNLean.PEPS.quantumDoubleOriginalTerm_eq_unblocked,
        TNLean.PEPS.quantumDoubleOriginalHamiltonian_eq_unblocked}

    \uses{def:peps_qd_original_hamiltonian}
    Let $U$ regroup each tile into its clockwise four-spin K register.
    If $\widetilde h_j$ and $\widetilde H$ are the placed K penalties and
    their sum, then on the entire physical space
    \begin{align}
        M_e(u)_{\sigma,\tau}
          &=\mathbf1_{\sigma=V_e(u)\tau},\\
        h_j&=U^\dagger\widetilde h_jU,&
        H&=U^\dagger\widetilde HU.
        \label{eq:peps_qd_original_blocking}
    \end{align}
\end{theorem}
\begin{proof}
    Regrouping~(\ref{eq:peps_qd_original_a}) gives exactly the local and
    hole K products, in their stated order. On a right B face the four
    affected spins receive $(R_u,L_u,R_u,L_u)$ in the order
    $(v,0),(v,1),(v+\mathrm e,2),(v+\mathrm e,3)$; on an upper B face
    they receive $(L_u,R_u,R_u,L_u)$ in the order
    $(v,0),(v,3),(v+\mathrm n,1),(v+\mathrm n,2)$.
    These are exactly the adjacent K-bond actions.
    The bijective regrouping sends every matrix entry to the corresponding
    K entry. Averaging with the same $|G|^{-1}$ and taking complements proves
    the B identity; the diagonal product-one indicators prove the A identity.
\end{proof}

\begin{theorem}[Original projectors and their complete common kernel]
    \label{thm:peps_qd_original_kernel}
    \lean{TNLean.PEPS.quantumDoubleOriginalATerm_mulVec_eq_zero_iff,
        TNLean.PEPS.quantumDoubleOriginalBTerm_mulVec_eq_zero_iff,
        TNLean.PEPS.quantumDoubleOriginalTerm_isStarProjection,
        TNLean.PEPS.quantumDoubleOriginalTerm_commute,
        TNLean.PEPS.quantumDoubleOriginalTerm_posSemidef,
        TNLean.PEPS.quantumDoubleOriginalHamiltonian_posSemidef,
        TNLean.PEPS.quantumDoubleOriginalHamiltonian_mulVec_eq_zero_iff,
        TNLean.PEPS.quantumDoubleOriginalHamiltonian_mulVec_eq_zero_iff_closureSpan}

    \uses{thm:peps_qd_original_blocking}
    All original penalties are commuting orthogonal projections and
    are positive semidefinite. Their sum is positive semidefinite. The
    coefficient conditions below are imposed for every $\sigma\in G^\Lambda$
    and $u\in G$:

    \begin{align}
        \ker H
          &=\bigcap_j\ker h_A^j\cap\bigcap_e\ker h_B^e,\\
        h_A^j\psi=0
          &\iff a_j(\sigma)\ne1\Longrightarrow\psi(\sigma)=0,\\
        h_B^e\psi=0
          &\iff\psi(V_e(u)\sigma)=\psi(\sigma).
    \end{align}
    Equivalently, $H\psi=0$ precisely when $U\psi$ belongs to the span
    of the actual commuting K closures.
\end{theorem}
\begin{proof}
    The A assertion follows from its diagonal entries. The finite-group
    average fixes exactly the invariant coefficient functions, proving
    the B assertion. The coefficientwise identities transport the
    full-space K commutators, orthogonal projections and complete kernel
    through the unitary $U$; positivity identifies the sum's kernel with
    the intersection of the individual kernels.
\end{proof}

\begin{theorem}[All original-spin ground sectors]
    \label{thm:peps_qd_original_sectors}
    \lean{TNLean.PEPS.quantumDoubleOriginalBlockingLinearEquiv,
        TNLean.PEPS.quantumDoubleOriginalBlockingLinearEquiv_apply,
        TNLean.PEPS.map_ker_quantumDoubleOriginalHamiltonian,
        TNLean.PEPS.quantumDoubleOriginalClosure,
        TNLean.PEPS.ker_quantumDoubleOriginalHamiltonian_eq_span_closures,
        TNLean.PEPS.finrank_ker_quantumDoubleOriginalHamiltonian}

    \uses{thm:peps_qd_original_kernel}
    Let $\Psi^K_{g,h}$ be the actual K closures and put
    $\Phi_{g,h}=U^\dagger\Psi^K_{g,h}$. Then
    \begin{align}
        \ker H&=\operatorname{span}_{\mathbb C}
          \{\Phi_{g,h}:gh=hg\},\\
        \dim_{\mathbb C}\ker H
          &=\#\bigl(\{(g,h):gh=hg\}/G\bigr),
    \end{align}
    where $G$ acts by simultaneous conjugation. No identity condition is imposed on either
    noncontractible torus holonomy.
\end{theorem}
\begin{proof}
    The linear regrouping maps the entire original kernel bijectively
    onto the K kernel. Its commuting-closure spanning theorem and
    commuting-pair conjugacy-class dimension formula therefore give
    both assertions.
\end{proof}

\begin{theorem}[The original orthogonal ground projector]
    \label{thm:peps_qd_original_ground_projector}
    \lean{TNLean.PEPS.quantumDoubleOriginalGroundProjector,
        TNLean.PEPS.quantumDoubleOriginalGroundProjector_eq_B_prod_A,
        TNLean.PEPS.quantumDoubleOriginalGroundProjector_isStarProjection,
        TNLean.PEPS.quantumDoubleOriginalGroundProjector_mulVec_eq_self_iff,
        TNLean.PEPS.range_quantumDoubleOriginalGroundProjector}

    \uses{thm:peps_qd_original_kernel}
    The correctly normalized product
    \begin{align}
        Q=\left(\prod_e P_B^e\right)\left(\prod_j(I-h_A^j)\right)
    \end{align}
    is an orthogonal projection with $\operatorname{ran}Q=\ker H$.
    Equivalently, $Q\psi=\psi$ if and only if $H\psi=0$.
\end{theorem}
\begin{proof}
    A product of commuting orthogonal projections is an orthogonal
    projection whose fixed space is the intersection of their fixed
    spaces. Here that intersection is the common kernel of the penalties.
\end{proof}

\begin{theorem}[Exact preparation from the group identity]
    \label{thm:peps_qd_original_preparation}
    \lean{TNLean.PEPS.quantumDoubleOriginalTerm_checkerboard,
        TNLean.PEPS.quantumDoubleOriginalHamiltonian_checkerboard,
        TNLean.PEPS.quantumDoubleOriginalCheckerboard_one_coefficient,
        TNLean.PEPS.quantumDoubleOriginalCheckerboard_ne_zero,
        TNLean.PEPS.quantumDoubleOriginalPhysicalColorConfig,
        TNLean.PEPS.quantumDoubleOriginalATerm_one,
        TNLean.PEPS.quantumDoubleOriginalGroundProjector_one_eq_B_product,
        TNLean.PEPS.quantumDoubleOriginalCheckerboard_eq_sum_colorBasis,
        TNLean.PEPS.card_quantumDoubleOriginalPhysicalColorings,
        TNLean.PEPS.quantumDoubleOriginalCheckerboard_eq_groundProjector_one,
        TNLean.PEPS.quantumDoubleOriginalCheckerboard_eq_B_product_one,
        TNLean.PEPS.quantumDoubleOriginalGroundProjector_one_ne_zero}

    \uses{thm:peps_qd_original_ground_projector,thm:peps_qd_original_blocking}
    Retain the elementary T tensor with its printed, unnormalized entries,
    and let $\Psi_T$ be its actual periodic contraction. Every original
    penalty annihilates $\Psi_T$, and
    \begin{align}
        \Psi_T
          &=|G|^{2wh}Q\ket{1,\ldots,1}
           =|G|^{2wh}\left(\prod_e P_B^e\right)\ket{1,\ldots,1},
        \label{eq:peps_qd_original_preparation}\\
        \Psi_T(1,\ldots,1)&=|G|\ne0.
    \end{align}
    Thus the projected product state is nonzero. The overall state scalar
    in~(\ref{eq:peps_qd_original_preparation}) does not change the required
    normalization of any individual projector.
\end{theorem}
\begin{proof}
    Physical blocking gives $U\Psi_T=\Psi_K$ with scalar one. The actual
    K contraction is the sum of the basis vectors produced by all shared
    edge colorings $c\in G^E$, so
    $\Psi_T=\sum_c\ket{\sigma(c)}$ with every stabilizer multiplicity retained.
    Every K penalty kills $\Psi_K$, hence $Q\Psi_T=\Psi_T$.
    Each column of $Q$ is a ground vector and is invariant under all bond
    color changes. Self-adjointness of $Q$ therefore gives
    $Q\ket{\sigma(c)}=Q\ket{1,\ldots,1}$ for every $c$.
    There are $|G|^{|E|}=|G|^{2wh}$ colorings, proving the first identity.
    The A projectors fix the all-identity input, so they can be omitted
    from its preparation. Finally the all-identity K coefficient is
    $|G|$, proving the nonzero assertion.
\end{proof}
